% Л2 -- assemblies, delivery models and memory. Slides.
%
% Companion to the spoken script in ../L02-assemblies-and-memory.md, which is
% the source of truth: the slides carry code, diagrams and short claims, and
% the delivery lives in \note{}. Slide text and speaker notes are in Russian
% because that is the content; everything else in this file is English.
%
% Build:  make          (latexmk -pdf; engine: pdflatex)
%         make notes    (same deck plus speaker notes on a second screen)
%
% Every terminal listing here is real output captured on the speaker's machine
% (dotnet SDK 10.0.401, ilspycmd 11.0.0). Paths are anonymised to /home/user
% and nothing else in them is touched.
%
% House rules: ../CLAUDE.md and ../../../CLAUDE.md. Theme: ../../../../theme/.

\documentclass[aspectratio=169,10pt,t]{beamer}

\usetheme{dotnet}
\usetikzlibrary{arrows.meta,positioning,calc,decorations.pathreplacing}

\dotnetlecture{Л2}

\ifdefined\shownotes
  \setbeameroption{show notes on second screen=right}
\fi

\title[Л2. Сборки и память]{Сборки, модели поставки и память}
\author{Дмитрий Курочкин}
\date{}

\begin{document}

% ------------------------------------------------------------------ title --
\begin{frame}[plain]
  \titlepagednc{Лекция 2}{Сборки, модели поставки и память}{}{Дмитрий Курочкин}
\end{frame}
\note{80 минут. Прошлую лекцию закончили на пяти файлах в bin -- сегодня
  вскрываем главный из них.}

% ============================================= 1. INSIDE AN ASSEMBLY =======
\section{1. Сборка изнутри}

\begin{frame}[plain]
  \sectionopener{1}{Сборка изнутри}{Что лежит внутри dll? К концу раздела вы
    сможете вскрыть любую библиотеку и прочитать её, как исходник.}
\end{frame}
\note{Бюджет: 20 минут. Сначала вопрос на повторение: что внутри dll. Кто
  скажет <<IL>> -- прав, но неполно.}

\begin{frame}{Сборка в разрезе}
  \slidetop
  \claim{В сборке \hi{три} вещи, и IL из них самая скучная.}

  \begin{center}
    \begin{tikzpicture}[font=\small,
      box/.style={draw=hairline,line width=0.8pt,fill=codebg,
                  minimum height=0.9cm,minimum width=3.3cm,align=center,
                  inner sep=0.35em},
      ar/.style={-{Straight Barb[length=1.8mm]},ink,line width=0.7pt}]
      \node[box] (man) {манифест};
      \node[box,below=0.45cm of man] (meta) {метаданные};
      \node[box,below=0.45cm of meta] (il) {IL};
      \node[anchor=west,muted] at ([xshift=0.6cm]man.east)
        {паспорт и список зависимостей};
      \node[anchor=west,muted] at ([xshift=0.6cm]meta.east)
        {таблицы всех типов и членов};
      \node[anchor=west,muted] at ([xshift=0.6cm]il.east)
        {тела методов};
      \draw[ar] (il.west) -- ++(-0.9,0) |- node[pos=0.25,left,
        font=\scriptsize\bfseries,text=accent,align=right]
        {ссылается\\на} (meta.west);
    \end{tikzpicture}
  \end{center}

  \slidegap
  \aside{IL нужен \lo{только рантайму} и только в момент выполнения.
    Метаданные нужны \hi{всем}.}
\end{frame}
\note{Метаданные читают компилятор, отладчик, рантайм, инструменты и ваш
  собственный код. Манифест -- как у корабля: что везём и что понадобится на
  другой стороне. Расширение dll -- наследство от Windows; на Linux файл открывает
  сам рантайм, и ему всё равно, как файл назван.}

\begin{frame}[fragile]{Программа знает о себе всё}
  \slidetop
\begin{lstlisting}[style=csharp,basicstyle=\ttfamily\scriptsize]
using System.Reflection;

Assembly assembly = typeof(Console).Assembly;
Console.WriteLine(assembly.FullName);

foreach (Type type in assembly.GetExportedTypes().Take(5))
{
    Console.WriteLine(type.FullName);
}
\end{lstlisting}
  \vspace{0.45em}
\begin{lstlisting}[style=term]
System.Console, Version=10.0.0.0, Culture=neutral, PublicKeyToken=b03f5f7f11d50a3a
System.Console
System.ConsoleCancelEventHandler
System.ConsoleCancelEventArgs
System.ConsoleColor
System.ConsoleSpecialKey
\end{lstlisting}
  \vspace{0.45em}
  \aside{Ни одного файла я не открывал и ни одного пути не указывал ---
    \hi{метаданные в действии}.}
\end{frame}
\note{Я просто спросил у работающей программы, из чего она состоит, и она
  ответила. Обратить внимание на первую строку вывода: там не просто имя, там
  четыре части через запятую. Про них -- следующий слайд.}

\begin{frame}{Идентичность сборки --- четыре части}
  \slidetop
  \claim{Она будет попадаться вам \hi{в сообщениях об ошибках}. Читать с
    листа.}

  \begin{center}
    {\ttfamily\small
     \hi{System.Console}, \hi{Version=10.0.0.0},
     {\color{muted}Culture=neutral, PublicKeyToken=b03f5f7f11d50a3a}}
  \end{center}

  \slidegap
  \begin{tabular}{@{}p{3.3cm}p{8.0cm}@{}}
    \textbf{\hi{имя}}        & обычно совпадает с именем файла, но не обязано \\[0.6em]
    \textbf{\hi{версия}}     & четыре числа; \lo{не} версия пакета NuGet \\[0.6em]
    \textbf{культура}        & для кода нейтральная; своя только у сборок с переводами \\[0.6em]
    \textbf{токен ключа}     & остаток строгих имён; рантайм подпись не проверяет \\
  \end{tabular}

  \slidegap
  \aside{В csproj одно свойство \texttt{Version}, а версий из него выводится
    \hi{три}: сборки, файла и информационная. В больших библиотеках их задают
    по отдельности.}
\end{frame}
\note{Strong name -- если встретите слова, теперь знаете, что это. Версия
  сборки -- то, на что смотрит рантайм.}

\begin{frame}{Две версии одной библиотеки}
  \slidetop
  \claim{В \texttt{bin} ляжет \hi{один} \texttt{Foo.dll}. Два файла с одним
    именем в одну папку не положишь.}

  \begin{center}
    \begin{tikzpicture}[font=\small,
      box/.style={draw=hairline,line width=0.8pt,fill=codebg,
                  minimum height=0.85cm,minimum width=2.6cm,align=center,
                  inner sep=0.35em},
      ar/.style={-{Straight Barb[length=1.8mm]},ink,line width=0.7pt}]
      \node[box] (a) {пакет А};
      \node[box,below=0.9cm of a] (b) {пакет Б};
      \node[box,right=4.2cm of a,yshift=-0.75cm,fill=paper,draw=accent]
        (foo) {\texttt{Foo.dll} 2.0};
      \draw[ar] (a.east) -- node[above,font=\scriptsize,text=muted,sloped]
        {Foo не ниже 1.2} (foo.west);
      \draw[ar] (b.east) -- node[below,font=\scriptsize,text=muted,sloped]
        {Foo не ниже 2.0} (foo.west);
      \node[muted,font=\scriptsize,below=0.35em] at (foo.south) {в bin};
    \end{tikzpicture}
  \end{center}

  \slidegap
  \hrulethin{0.42}
  \rulegap
  {\color{alt}\large\bfseries Метод не найден --- при первом вызове, уже
   в проде.}
\end{frame}
\note{Голосование в три руки: 2.0, 1.2 или ошибка сборки. NuGet выберет одну
  версию на всё приложение -- 2.0, она удовлетворяет обоих. Ломается, только
  если из 2.0 убрали метод, который зовёт пакет А. Сборка при этом проходит:
  компилятор проверял ваш код, а не код пакета А. Запуск тоже: сборки
  загружаются лениво.}

\begin{frame}{Как рантайм находит сборки}
  \slidetop
  \claim{Два места: \hi{папка приложения} и папка фреймворка. Больше нигде.}

  \begin{center}
    \begin{tikzpicture}[font=\small,
      box/.style={draw=hairline,line width=0.8pt,fill=codebg,
                  minimum height=1.15cm,minimum width=2.2cm,align=center,
                  inner sep=0.28em,font=\scriptsize},
      ar/.style={-{Straight Barb[length=1.8mm]},ink,line width=0.7pt}]
      \node[box] (l) {лаунчер};
      \node[box,right=0.45cm of l] (rc) {\texttt{runtimeconfig}\\\texttt{.json}};
      \node[box,right=0.45cm of rc] (rt) {рантайм\\нужной версии};
      \node[box,right=0.45cm of rt] (dj) {\texttt{deps.json}};
      \node[box,right=0.45cm of dj] (list) {перечень сборок:\\своя папка\\и папка фреймворка};
      \draw[ar] (l) -- (rc);
      \draw[ar] (rc) -- (rt);
      \draw[ar] (rt) -- (dj);
      \draw[ar] (dj) -- (list);
      \node[accent,font=\scriptsize\bfseries,below=0.5em] at (list.south)
        {и только потом -- загрузка};
    \end{tikzpicture}
  \end{center}

  \slidegap
  \aside{Загрузка \hi{ленивая}: сборка читается при первом обращении к типу.
    Старт от этого быстрее, но пропавшая сборка обнаружится
    \lo{не при старте, а посреди работы}.}
\end{frame}
\note{Никаких системных папок и общего хранилища на машину, как было в
  .NET Framework. Точка расширения называется контекст загрузки сборок --
  через неё делают плагины; слово запомнить.}

\begin{frame}[fragile]{Вскрыть сборку --- это одна команда}
  \slidetop
\begin{lstlisting}[style=term]
$ ilspycmd bin/Release/net10.0/HelloWorld.dll
using System;
...
[assembly: TargetFramework(".NETCoreApp,Version=v10.0", FrameworkDisplayName = ".NET 10.0")]
[assembly: AssemblyFileVersion("1.0.0.0")]
[assembly: AssemblyInformationalVersion("1.0.0")]
[assembly: AssemblyVersion("1.0.0.0")]
[CompilerGenerated]
internal class Program
{
    private static void Main(string[] args)
    {
        Console.WriteLine("Hello, World!");
    }
}
\end{lstlisting}
  \vspace{0.7em}
  \aside{Одна строка в \texttt{Program.cs} --- а в файле \hi{класс и Main},
    которых вы не писали, и \hi{три версии} из одного свойства. Класс
    \lo{internal}, метод \lo{private} --- модификаторы доступа это правила
    для компилятора, а не защита.}
\end{frame}
\note{То же самое работает на любой чужой dll прямо из кэша пакетов: приватные
  поля, internal-типы, строковые литералы. Два вывода: секрет в коде не
  спрячешь, и заглядывать внутрь пакета -- это не хакерство, это чтение.
  В шестой лекции так же посмотрим на асинхронный метод.}

\begin{frame}[plain]
  \keyline{Сборка --- не чёрный ящик, а читаемая структура:\\ паспорт,
    оглавление и код. Кто угодно может её открыть, и вы тоже.}
\end{frame}
\note{И запомните на потом: через полчаса увидим модели поставки, которые эту
  самоописанность разрушают.}

% ==================================================== 2. NUGET =============
\section{2. NuGet}

\begin{frame}[plain]
  \sectionopener{2}{NuGet}{Проект $\neq$ сборка $\neq$ пакет. Три разные
    вещи, три разные единицы.}
\end{frame}
\note{Бюджет: 15 минут. История с двумя версиями Foo на этом не
  заканчивается.}

\begin{frame}{NuGet --- это три вещи}
  \slidetop
  \claim{Одним словом называют \hi{три} разные вещи, и путаница стоит
    времени.}

  \begin{tabular}{@{}p{5.4cm}p{6.0cm}@{}}
    \textbf{\hi{формат}} \texttt{.nupkg}  & во что упаковано \\[0.8em]
    \textbf{\hi{клиент}} в \texttt{dotnet} CLI & кто качает и подставляет
                                                 в сборку \\[0.8em]
    \textbf{\hi{хранилище}} \texttt{nuget.org} или свой сервер
                                          & откуда \\
  \end{tabular}

  \slidegap
  \hrulethin{0.42}
  \rulegap
  \aside{<<Возьми из NuGet>> --- это \hi{третье}. <<NuGet ругается на
    версии>> --- \lo{второе}.}
\end{frame}
\note{Дальше разберём все три. Сначала формат.}

\begin{frame}[fragile]{Пакет --- это архив}
  \slidetop
  \begin{minipage}[t]{0.55\textwidth}
\begin{lstlisting}[style=term]
$ unzip -l Pure.Primitives.3.6.5.nupkg
  Length      Name
---------  ---------------------------------
     1558  Pure.Primitives.nuspec
    74752  lib/net10.0/Pure.Primitives.dll
    75264  lib/net7.0/Pure.Primitives.dll
    74752  lib/net8.0/Pure.Primitives.dll
    74752  lib/net9.0/Pure.Primitives.dll
\end{lstlisting}
  \end{minipage}\hfill
  \begin{minipage}[t]{0.41\textwidth}
    \colhead{accent}{внутри всегда три вещи}\par\vspace{0.8em}
    \begin{cbul}{accent}
      \item сборки --- по одной на каждую целевую платформу
      \item метаданные: имя, версия, автор, лицензия
      \item список зависимостей, свой для каждой платформы
    \end{cbul}
    \vspace{0.8em}
    \colhead{alt}{бывает и четвёртое}\par\vspace{0.5em}
    \aside{анализаторы, встраиваемые файлы, правила для MSBuild}
  \end{minipage}

  \vspace{0.6em}
  \aside{Обычный zip с другим расширением. Пакет версии \hi{3.6.5}, сборка
    внутри --- \lo{1.0.0.0}.}
\end{frame}
\note{Версия пакета и версия сборки внутри него -- разные числа, и это не
  ошибка. restore смотрит на ваш TargetFramework и берёт подходящую папку.
  Проект -- единица сборки. Сборка -- единица загрузки и идентичности.
  Пакет -- единица распространения: внутри может быть пять сборок, а может
  не быть ни одной. Авторы держат версию сборки стабильной в пределах
  мажорной, чтобы не ломать ссылки.}

\begin{frame}{Версия в PackageReference --- не <<ровно эта>>}
  \slidetop
  \claim{Голая версия --- это \hi{нижняя граница} диапазона.}

  \begin{tabular}{@{}p{4.2cm}p{7.2cm}@{}}
    \texttt{\hi{13.0.3}}       & не ниже 13.0.3 --- \hi{так пишут обычно} \\[0.8em]
    \texttt{\lo{[13.0.3]}}     & ровно 13.0.3, и никакая другая \\[0.8em]
    \texttt{\lo{[13.0,14.0)}}  & диапазон: 13.0 включая, 14.0 исключая \\
  \end{tabular}

  \slidegap
  \hrulethin{0.42}
  \rulegap
  \aside{\texttt{restore} берёт \hi{минимальную подходящую}. Квадратная
    скобка включает границу, круглая исключает --- синтаксис из математики,
    читается так же. Побеждает \hi{ближайший}: прямая ссылка из вашего
    csproj важнее того, что просит транзитивная зависимость.}
\end{frame}
\note{Отсюда стандартный приём: уязвимость в транзитивном пакете лечится
  прямой ссылкой с нужной версией, не дожидаясь автора. Понижение версии
  restore предупредит -- предупреждение стоит читать.}

\begin{frame}[fragile]{Свой пакет: четыре свойства и две команды}
  \slidetop
\begin{lstlisting}[style=xml]
<PropertyGroup>
  <PackageId>Bsuir.Course.Shared</PackageId>
  <Version>1.0.0</Version>
  <Authors>Your Name</Authors>
  <Description>Shared domain types for the course service.</Description>
</PropertyGroup>
\end{lstlisting}
  \vspace{0.8em}
\begin{lstlisting}[style=bash,basicstyle=\ttfamily\scriptsize]
dotnet pack -c Release
dotnet nuget push bin/Release/Bsuir.Course.Shared.1.0.0.nupkg --source <feed> --api-key <key>
\end{lstlisting}
  \vspace{0.8em}
  \aside{Пакет с одной и той же версией нельзя опубликовать дважды, и это
    правильно: версия --- это обещание, что содержимое под этим номером
    \hi{не меняется никогда}. Нашли ошибку --- выпускаете патч, а не
    перезаливаете старый.}
\end{frame}
\note{В компании хранилище почти всегда свой сервер: разница только в адресе
  источника. Ломаете совместимость -- поднимайте мажорную. Ещё в пакет
  кладут лицензию, ссылку на репозиторий, описание и символы отдельным
  пакетом.}

\begin{frame}{Четыре вопроса перед \texttt{add package}}
  \slidetop
  \claim{Зависимость --- обязательство \hi{на годы}, а не на вечер.}

  \begin{enumerate}\itemsep0.6em
    \item \textbf{Проект живой?} Когда последний коммит, отвечают ли на
          вопросы, сколько людей поддерживают. Не трогали три года --- это
          не \lo{<<стабильная>>}.
    \item \textbf{Сколько тянет за собой?} Посмотрите транзитивные. Иногда
          функцию проще написать за час, чем годами обслуживать дерево,
          которое пришло за ней.
    \item \textbf{Какая лицензия?} Есть те, что требуют открыть и ваш код.
          Есть коммерческие. Юристы спросят --- лучше раньше.
    \item \textbf{Выкинете за вечер?} За своим интерфейсом --- да. Если типы
          расползлись по всему коду --- \hi{вы с ней навсегда}.
  \end{enumerate}

  \slidegap
  \aside{\texttt{restore} сам предупредит об известной уязвимости ---
    не отключайте. И смотрите на имя: пакеты, отличающиеся одной буквой,
    существуют.}
\end{frame}
\note{Чужой код будет выполняться в вашем процессе, с вашими правами, с
  доступом к вашим данным. Одна строка -- и это уже так.}

\begin{frame}[plain]
  \keyline{Пакет --- это единица распространения, а не сборка.\\ Каждая
    строка \texttt{PackageReference} --- обязательство.}
\end{frame}
\note{Понимать его стоит до того, как взять, а не после.}

% =========================================== 3. DELIVERY MODELS ============
\section{3. Модели поставки}

\begin{frame}[plain]
  \sectionopener{3}{Модели поставки}{%
    \normalcolor
    \begin{tikzpicture}[font=\scriptsize,
      rung/.style={draw=hairline,line width=0.8pt,fill=codebg,
                   minimum height=0.62cm,minimum width=2.45cm,align=center,
                   inner sep=0.2em}]
      \foreach \i/\n in {0/framework-dependent, 1/self-contained,
                         2/ReadyToRun, 3/trimming, 4/NativeAOT} {
        \node[rung] at (\i*2.15,\i*0.52) {\texttt{\n}};
      }
      \node[muted,anchor=west] at (-1.15,2.08) {};
    \end{tikzpicture}\par
    \vspace{0.9em}
    Пять ступеней. На каждой вы что-то получаете и за что-то платите ---
    лучшего варианта нет, есть подходящий под задачу.}
\end{frame}
\note{Бюджет: 20 минут. Вопрос в зал: что нужно на машине пользователя, чтобы
  запустить HelloWorld из прошлой лекции? Правы те, кто скажет <<нужен
  .NET>>, -- но правы не всегда.}

\begin{frame}{Обрезка: что видно статически}
  \slidetop
  \claim{Анализатор идёт от точки входа по вызовам. Что не достижимо ---
    \lo{мусор}.}

  \begin{center}
    \begin{tikzpicture}[font=\scriptsize,
      n/.style={draw=hairline,line width=0.8pt,fill=codebg,
                minimum height=0.7cm,minimum width=2.1cm,align=center,
                inner sep=0.25em},
      ar/.style={-{Straight Barb[length=1.8mm]},ink,line width=0.7pt}]
      \node[n,draw=accent] (e) {точка входа};
      \node[n,above right=0.35cm and 1.5cm of e] (a) {тип A};
      \node[n,below right=0.35cm and 1.5cm of e] (r)
        {найти тип\\по строке};
      \node[n,right=1.5cm of a] (b) {тип B};
      \node[n,right=1.5cm of r,draw=alt,dashed,text=alt] (x) {тип C};
      \draw[ar] (e) -- (a);
      \draw[ar] (e) -- (r);
      \draw[ar] (a) -- (b);
      \draw[-{Straight Barb[length=1.8mm]},alt,line width=0.7pt,dashed]
        (r) -- node[above,font=\scriptsize\bfseries,text=alt]{?} (x);
      \node[alt,font=\scriptsize\bfseries,below=0.3em] at (x.south)
        {выкинут};
    \end{tikzpicture}
  \end{center}

  \slidegap
  \aside{Какое имя будет в строке, известно только во время выполнения.
    Поэтому ошибка \lo{не при сборке, а при выполнении} --- и только
    в обрезанной версии: в обычной сборке всё работало.}
\end{frame}
\note{Рефлексией пользуется гораздо больше кода, чем вы думаете: поиск
  контроллеров по соглашению (Л9), сериализатор (Л9), ORM (Л14), внедрение
  зависимостей (Л10). Если соберётесь применить её сами -- сначала
  посчитайте, что отдаёте: обрезку, AOT, проверки компилятора и
  предсказуемость.}

\begin{frame}[fragile]{Пять ступеней --- пять команд}
  \slidetop
\begin{lstlisting}[style=bash,basicstyle=\ttfamily\scriptsize]
dotnet publish -c Release
dotnet publish -c Release -r linux-x64 --self-contained
dotnet publish -c Release -r linux-x64 -p:PublishReadyToRun=true
dotnet publish -c Release -r linux-x64 --self-contained -p:PublishTrimmed=true
dotnet publish -c Release -r linux-x64 -p:PublishAot=true
\end{lstlisting}
  \vspace{1.0em}
  \aside{Первая \hi{без RID}, потому что IL универсален. Остальные с RID,
    потому что везде появляется нативный код. Обрезка требует
    самодостаточности, иначе нечего обрезать; AOT подразумевает её сам.}

  \slidegap
  \aside{На практике эти свойства пишут \hi{не в командной строке, а в
    csproj} --- тогда анализаторы включаются уже на обычной сборке и
    предупреждают заранее, а не в момент публикации.}
\end{frame}
\note{Ровно это вы сделаете в лабораторной: опубликуете одно и то же
  приложение тремя способами и увидите разницу своими глазами.}

\begin{frame}{Пять моделей в одной таблице}
  \slidetop
  {\scriptsize
  \begin{tabular}{@{}p{3.3cm}p{1.9cm}p{1.9cm}p{1.9cm}p{3.1cm}@{}}
    & \colhead{accent}{нужен .NET} & \colhead{accent}{размер}
    & \colhead{accent}{старт} & \colhead{alt}{перестаёт работать} \\[0.9em]
    \texttt{framework-dependent} & да  & минимальный & прогрев JIT
      & ничего \\[0.7em]
    \texttt{self-contained}      & нет & большой     & прогрев JIT
      & ничего \\[0.7em]
    \texttt{ReadyToRun}          & по выбору & ещё больше & \hi{быстрый}
      & ничего \\[0.7em]
    \texttt{trimming}            & нет & средний     & прогрев JIT
      & \lo{рефлексия на невидимое статически} \\[0.7em]
    \texttt{NativeAOT}           & нет & \hi{малый}, один файл
      & \hi{мгновенный}
      & \lo{динамическая загрузка, генерация кода, часть рефлексии} \\
  \end{tabular}}

  \slidegap
  \aside{Платите вы \lo{динамикой} --- которой в нормальном прикладном коде
    быть и не должно. Тогда цена пятой ступени для вас почти нулевая.}
\end{frame}
\note{Патч безопасности: на первой ступени его ставят один раз на машину и
  все приложения исправлены; на остальных это ваша пересборка. Сто
  микросервисов со своим рантаймом внутри -- сто пересборок на каждый патч.}

\begin{frame}{Какой рантайм возьмёт приложение}
  \slidetop
  \claim{Версия в \texttt{runtimeconfig} --- это \hi{минимум}, с которого
    начинается поиск.}

  \begin{center}
    \begin{tikzpicture}[font=\small,
      box/.style={draw=hairline,line width=0.8pt,fill=codebg,
                  minimum height=0.85cm,minimum width=3.1cm,align=center,
                  inner sep=0.35em},
      ar/.style={-{Straight Barb[length=1.8mm]},line width=0.8pt}]
      \node[box] (rc) {\texttt{runtimeconfig}\\просит \texttt{10.0.0}};
      \node[box,right=3.6cm of rc,yshift=0.85cm,draw=accent] (a)
        {\texttt{10.0.5}};
      \node[box,right=3.6cm of rc,yshift=-0.85cm,draw=hairline] (b)
        {\texttt{11.0.1}};
      \node[muted,font=\scriptsize,above=0.3em] at (a.north)
        {установлено на машине};
      \draw[ar,accent] (rc.east) -- node[above,font=\scriptsize\bfseries,
        text=accent,sloped]{патч посвежее} (a.west);
      \draw[ar,alt,dashed] (rc.east) -- node[pos=0.5,fill=paper,text=alt,
        font=\Large\bfseries,inner sep=1.5pt]{$\times$} (b.west);
      \node[alt,font=\scriptsize\bfseries,below=0.3em] at (b.south)
        {через мажорную границу --- по умолчанию нет};
    \end{tikzpicture}
  \end{center}

  \slidegap
  \aside{Именно поэтому патчи безопасности \hi{доезжают до приложений сами}.
    Мажорная версия может ломать совместимость --- рантайм не решает за вас.
    Политика меняется \lo{снаружи, без пересборки}.}
\end{frame}
\note{Когда это ваша проблема: на сервере обновили рантайм, и что-то
  изменилось, хотя вы ничего не выкладывали. Выводите версию рантайма в лог
  при старте -- через полгода на разборе инцидента эта строка сэкономит час.}

% ================================================= 4. MEMORY AND GC ========
\section{4. Память и GC}

\begin{frame}[plain]
  \sectionopener{4}{Память и GC}{За двадцать минут я не научу вас управлять
    памятью. Я дам словарь, чтобы читать серьёзные книги.}
\end{frame}
\note{Бюджет: 20 минут. Вопрос в зал: число в локальной переменной и список
  через new -- где что живёт? Простая модель верна, и мы её сломаем.}

\begin{frame}{Стек и куча}
  \slidetop
  \claim{Значение живёт \hi{там, где объявлено}, а не <<на стеке>>.}

  \begin{center}
    \begin{tikzpicture}[font=\scriptsize,
      fr/.style={draw=hairline,line width=0.8pt,fill=codebg,
                 minimum height=0.62cm,minimum width=3.0cm,align=center},
      ob/.style={draw=hairline,line width=0.8pt,fill=codebg,
                 minimum size=0.62cm,align=center},
      ar/.style={-{Straight Barb[length=1.8mm]},line width=0.7pt}]
      \node[muted] at (0,2.25) {СТЕК ПОТОКА};
      \node[fr] (f3) at (0,1.55) {кадр \texttt{Parse}};
      \node[fr] (f2) at (0,0.88) {кадр \texttt{Read}};
      \node[fr] (f1) at (0,0.21) {кадр \texttt{Main}};
      \draw[ar,accent] (-2.25,1.55) -- node[left,text=accent,
        font=\scriptsize\bfseries,align=right]{указатель\\сдвинулся} ++(0.6,0);
      \node[muted] at (6.3,2.25) {УПРАВЛЯЕМАЯ КУЧА};
      \draw[hairline,line width=0.8pt] (4.0,-0.25) rectangle (8.6,1.9);
      \node[ob] (o1) at (4.75,1.35) {};
      \node[ob] (o2) at (5.9,0.85) {};
      \node[ob] (o3) at (7.3,1.4) {};
      \node[ob] (o4) at (7.9,0.4) {};
      \draw[ar,ink] (f3.east) -- node[above,font=\scriptsize,text=muted]
        {ссылка} (o1.west);
      \draw[ar,ink] (o1) -- (o3);
      \draw[ar,ink] (o2) -- (o4);
    \end{tikzpicture}
  \end{center}

  \slidegap
  \aside{Выделить на стеке --- \hi{сдвинуть один указатель}; освободить ---
    сдвинуть обратно. Но память живёт ровно столько, сколько живёт вызов.
    Стек невелик и фиксирован: его переполнение --- \lo{единственная ошибка
    в .NET, которую нельзя поймать}, процесс завершается без исключения.}
\end{frame}
\note{Поле int внутри класса? Внутри объекта, в куче. Структура локально? В
  кадре. Массив из тысячи int? Сам массив в куче, все значения внутри него
  подряд. Число, захваченное лямбдой? Компилятор перенесёт его в объект в
  куче. И упаковка, boxing: каждая -- аллокация, каждая распаковка --
  проверка и копирование.}

\begin{frame}{Сборщик мусора --- это четыре роли}
  \slidetop
  \begin{tabular}{@{}p{3.4cm}p{8.0cm}@{}}
    \textbf{\hi{аллокатор}}   & выдать место под новый объект \\[0.7em]
    \textbf{\hi{модификатор}} & ваш собственный код: он меняет граф объектов,
                                и за ним надо успевать \\[0.7em]
    \textbf{\hi{коллектор}}   & найти мусор по достижимости \hi{от корней} \\[0.7em]
    \textbf{компактор}        & уплотнить, чтобы дальше было куда класть \\
  \end{tabular}

  \slidegap
  \hrulethin{0.42}
  \rulegap
  \eyebrow{вопрос в зал}\par
  \vspace{0.5em}
  {\Large Является ли \texttt{new} частью сборщика мусора?}
\end{frame}
\note{Правы те, кто скажет <<да>>. new -- это сдвиг указателя внутри куска,
  который рантайм уже забрал у системы; часто быстрее, чем malloc. Но у
  сборщика есть бюджет, и когда он исчерпан, очередной new запускает сборку.
  Вы платите за аллокацию не в момент new, а позже, оптом. Отсюда термин:
  давление на сборщик. Циклические ссылки не проблема: считается
  достижимость от корней, а не ссылки.}

\begin{frame}{Компактация: автобус}
  \slidetop
  \claim{Свободно двести мегабайт, массив на сто \lo{не помещается} ---
    это фрагментация.}

  \begin{center}
    \begin{tikzpicture}[font=\scriptsize,x=0.52cm,
      seat/.style={draw=hairline,line width=0.6pt,minimum size=0.42cm,
                   inner sep=0pt},
      full/.style={seat,fill=accent!25!paper},
      free/.style={seat,fill=paper}]
      \node[muted,anchor=east] at (-0.6,0) {ДО};
      \foreach \i/\s in {0/full,1/free,2/full,3/free,4/free,5/full,6/free,
                         7/full,8/free,9/free,10/full,11/free,12/free,
                         13/full,14/free,15/free} {
        \node[\s] at (\i,0) {};
      }
      \node[muted,anchor=east] at (-0.6,-1.15) {ПОСЛЕ};
      \foreach \i in {0,...,5} { \node[full] at (\i,-1.15) {}; }
      \foreach \i in {6,...,15} { \node[free] at (\i,-1.15) {}; }
      \draw[decorate,decoration={brace,amplitude=4pt,mirror},hairline]
        (5.7,-1.45) -- (15.3,-1.45);
      \node[accent,font=\scriptsize\bfseries] at (10.5,-2.0)
        {свободно одним куском};
    \end{tikzpicture}
  \end{center}

  \slidegap
  \aside{Объекты двигаются, значит адрес меняется. Поэтому вам и не дают
    сырой указатель: ссылка --- это не адрес, а то, что \hi{рантайм умеет
    обновить}. И поэтому же на часть сборки потоки \lo{останавливаются}.}
\end{frame}
\note{Безопасность памяти и уплотнение -- одна история с двух сторон.
  Водитель просит всех пересесть вперёд. Все сдвинулись, дыры исчезли,
  сзади сплошной ряд. В C вы получаете указатель, которым можно
  промахнуться, и память, которую надо освобождать руками. Здесь эти ошибки
  невозможны by design -- пока вы не написали unsafe.}

\begin{frame}{Фоновая сборка}
  \slidetop
  \claim{Полная сборка кучи в гигабайты \hi{не} означает, что сервис молчит
    секунду.}

  \begin{center}
    \begin{tikzpicture}[font=\scriptsize,x=0.86cm,y=0.18cm]
      \fill[alt!22!paper] (2.6,-0.6) rectangle (2.95,11.6);
      \fill[alt!22!paper] (8.3,-0.6) rectangle (8.65,11.6);
      \foreach \i in {0,...,9} {
        \draw[accent,line width=1.4pt] (0,\i) -- (11,\i);
      }
      \foreach \i in {10,11} {
        \draw[alt,line width=1.4pt] (0,\i) -- (11,\i);
      }
      \node[accent,anchor=west,font=\scriptsize\bfseries] at (11.3,4.5)
        {10 потоков: ваш код};
      \node[alt,anchor=west,font=\scriptsize\bfseries] at (11.3,10.5)
        {2: сборка gen 2};
      \node[alt,anchor=south,font=\scriptsize\bfseries] at (2.78,12.2)
        {общая остановка};
      \draw[-{Straight Barb[length=1.8mm]},muted,line width=0.6pt]
        (0,-2.6) -- node[below,font=\scriptsize,text=muted]{время} (11,-2.6);
    \end{tikzpicture}
  \end{center}

  \slidegap
  \aside{Останавливаются все --- но это \hi{две узкие полоски}, а не весь
    прямоугольник. Цена видна там же: \lo{два ядра из двенадцати} заняты не
    вашей работой.}
\end{frame}
\note{Самую дорогую часть -- сборку второго поколения -- сборщик делает в
  основном параллельно с вашими потоками. Эти два ядра и есть то
  процессорное время, которое вы увидите в счётчике <<процент времени в
  сборке мусора>>.}

\begin{frame}{Словарь}
  \slidetop
  \begin{tabular}{@{}p{4.0cm}p{7.4cm}@{}}
    \textbf{\hi{поколения 0/1/2}} & большинство объектов умирает молодыми:
      молодых проверяют часто и быстро, старых --- редко \\[0.65em]
    \textbf{куча больших объектов} & от \textasciitilde 85 КБ, в основном
      массивы; по умолчанию \lo{не уплотняется} \\[0.65em]
    \textbf{Workstation и Server} & бережёт отзывчивость \emph{против}
      по куче и потоку на ядро; для веба SDK включает второй сам \\[0.65em]
    \textbf{background GC} & дорогая часть сборки идёт параллельно вашему
      коду \\[0.65em]
    \textbf{\hi{escape analysis}} & JIT доказал, что ссылка не убегает из
      метода, --- объект ляжет на стек \\[0.65em]
    \textbf{\hi{Span}} & окно в чужую память: начало и длина, \hi{без
      копирования}; живёт только на стеке \\
  \end{tabular}
\end{frame}
\note{Самый дорогой объект -- не старый и не молодой, а тот, что прожил ровно
  столько, чтобы его дважды повысили. В сервисе так себя ведут объекты на
  один долгий запрос. И: .NET резервирует память с запасом и отдавать не
  спешит -- <<сервис съел четыре гигабайта>> само по себе не диагноз.
  Смотреть надо на размер управляемой кучи и растёт ли он от сборки к
  сборке.}

\begin{frame}[fragile]{Три функции, один результат}
  \slidetop
\begin{lstlisting}[style=csharp,basicstyle=\ttfamily\scriptsize]
static int YearOld(string line)
{
    string date = line.Substring(0, line.IndexOf(';'));
    return int.Parse(date.Substring(0, 4));
}
static int YearNew(string line)
{
    ReadOnlySpan<char> date = line.AsSpan(0, line.IndexOf(';'));
    return int.Parse(date.Slice(0, 4));
}
static int YearNewer(string line)
{
    ReadOnlySpan<char> date = line.AsSpan(0, line.IndexOf(';'));
    return int.Parse(date[..4]);
}
\end{lstlisting}
  \vspace{0.7em}
  \aside{Первая: \lo{две промежуточные строки в куче}, каждая под бюджет
    сборщика. Вторая: \hi{два окна, ноль аллокаций}. Третья --- та же
    вторая, синтаксисом диапазонов: \texttt{date[..4]} и
    \texttt{date.Slice(0, 4)} буквально одно и то же. Читаемость почти не
    пострадала.}
\end{frame}
\note{Span окупается там, где это делается на каждый запрос миллион раз в
  день: разбор текста, буферы. В коде, который выполняется раз в секунду,
  разница не стоит ни минуты вашего времени. Не боритесь с аллокациями
  вслепую: сначала измерить, потом оптимизировать, и измерять в Release --
  escape analysis это оптимизация JIT.}

\begin{frame}{Куда копать}
  \slidetop
  \claim{Этой теме посвящены \hi{целые книги}. Вот две двери.}

  \begin{tabular}{@{}p{5.6cm}p{5.8cm}@{}}
    \textbf{\hi{Pro .NET Memory Management}} & Конрад Кокоса. Толстая,
      подробная, лучшее, что есть по теме \\[0.9em]
    \textbf{\hi{Book of the Runtime}} & раздел документации в репозитории
      \texttt{dotnet/runtime}, написанный теми, кто сборщик и пишет \\
  \end{tabular}

  \slidegap
  \hrulethin{0.42}
  \rulegap
  \aside{Мы не трогали \texttt{write barriers}, \texttt{card tables},
    \texttt{pinning} и настройки бюджета аллокаций. Слова названы, чтобы вы
    узнали их при встрече. \lo{Глубина за вами.}}
\end{frame}
\note{Memory -- старший брат Span, который можно сохранить; про него в Л5
  вместе с коллекциями.}

\begin{frame}[plain]
  \keyline{Память в .NET безопасна ровно потому,\\ что вы ею не управляете.
    Цена --- сборщик: знать его нужно настолько, чтобы ему не мешать.}
\end{frame}
\note{Всё остальное в разделе -- следствия.}

% =========================================== 5. WHEN THINGS GO BAD =========
\section{5. Чем смотреть}

\begin{frame}{Чем смотреть, когда стало плохо}
  \slidetop
  \claim{Посмотреть на \hi{работающий процесс}, не трогая его.}

  \begin{tabular}{@{}p{3.6cm}p{7.8cm}@{}}
    \textbf{\hi{dotnet-counters}} & что именно плохо: живые числа раз в
      секунду --- куча, поколения, процент времени в сборке, исключения,
      потоки \\[0.8em]
    \textbf{\hi{dotnet-trace}} & где горит процессор: запись за отрезок
      времени, какие методы выполнялись и сколько \\[0.8em]
    \textbf{\hi{dotnet-dump}} & кто держит память: слепок кучи, какие типы
      её занимают и почему сборщик их не собрал \\
  \end{tabular}

  \slidegap
  \hrulethin{0.42}
  \rulegap
  \aside{Все три подключаются к работающему процессу по идентификатору:
    \hi{пересобирать и перезапускать не нужно}. Счётчики почти бесплатны,
    трасса стоит заметно, \lo{дамп --- это остановка}.}
\end{frame}
\note{Ставятся все три через dotnet tool install -g -- тем самым
  механизмом из прошлой лекции. Честный ответ <<добавлю логов и
  перевыложу>> стоит часов, и после
  перевыкладки проблема может не повториться. Все три работают в контейнере
  и на сервере: это не учебные игрушки, это то, чем чинят прод. В Л12 те же
  счётчики будем экспортировать в мониторинг, чтобы они рисовались всегда,
  а не когда вы вспомнили подключиться.}

\end{document}
